\documentclass[12pt,a4paper]{article}
\usepackage{geometry}
\usepackage[english]{babel}
\usepackage[utf8]{inputenc}
\usepackage{amsmath}
\usepackage{amssymb}
\usepackage{graphicx}
\usepackage{hyperref}
\usepackage{xcolor}
\usepackage{todonotes}
\setuptodonotes{inline}

\begin{document}
\thispagestyle{empty}
\begin{center}
Rheinisch-Westf\"alische Technische Hochschule Aachen \\
Alessandro Berti\\[6ex]

\LARGE
	\textbf{Phase Review: Testing and Assessment} \\[6ex]
        \centering{Group 4:}\\
	\textit{Richard Schuppener, Cem Colak} \\[6ex]
\end{center}

\newpage

\tableofcontents
\newpage

\section{Summary}
In this phase, we focused on gaining actual insights about the central question "Does conformance information derived from a process model improve the prediction of future deviations and compliance risks?". 
\subsection{Incorporated Improvement Suggestions}
We considered Alessandro's improvement suggestions and packaged the whole software artifact as an installable python project, and cleaned out obsolete CLI output artifacts and the HTML-based documentation from our Code folder in the repository. Cleaning the repository was done by simply editing the \texttt{gitignore} file.
The software artifact can be installed by, from the Git repository root directory, running:\\
\texttt{pip install -e .}\\ 
if the software is to be installed in editable mode (this means that changes to the .py files directly take effect on the installed package).
If no further development is desired, the software artifact can also be installed in standard mode via:\\
\texttt{pip install .}\\ 
(from the Git repository root directory). Alternatively, the CLI tool can also be installed without cloning the Git repository by running the following command:\\
\texttt{pip install git+https://git.rwth-aachen.de/schuppenerr/conformance-predictions.git}\\


\section{Sprint review}

\subsection{Richard Schuppener}
This last phase of the lab was for me mainly defined by it's rewarding feeling that we finally completed the software and could show our results. Packaging the software artifact gave me the feeling that it was now complete.
Evaluating our results was a tedious task and only possible in this short time frame due to our proactive measures in the second and third sprint, where we implemented a multitude of visualizations, which we could now facilitate to demonstrate our software's performance.
Again, I am very content with the communicaton and collaboration with Cem. 
\subsection{Cem Colak}
In this phase of the lab I mainly focused on building an evaluation case which took a lot of time and revealed some heavy bottlenecks of my own device and therefore also for real life applications despite possible optimizations. The \emph{future conformance deviations} prediction target took too much time, such that I can't meaningfully evaluate its perfomance and reliability in terms of the outputs it produces. Therefore I switched to remaining time as a prediction target and delivered a hopefully meaningful evaluation. Despite this trouble, I did some improvements on the remaining time evaluations which now provide reliable and straight-forward outputs. Me and Richard initially talked about how our output format for this project should look like and discussed some detail before we worked independently on our parts of the evaluation. I really enjoyed the whole lab with him together and also this final stage of our lab.
\end{document}